\documentclass[10pt,a4paper]{amsart}
\usepackage[margin=19mm]{geometry}
\usepackage{amsmath,amssymb,mathtools}
\usepackage[scr=rsfs,cal=euler]{mathalfa}
\usepackage{xfrac}
\usepackage[unicode,hidelinks,bookmarks=false]{hyperref}
\newcommand{\ie}{i.e.\ }
\newcommand{\cf}{cf.\ }
\newcommand{\resp}{respectively\ }
\newcommand{\Subsection}{\subsection}
\providecommand{\nobreakdash}{\nobreak\hbox{-}}
\setlength{\emergencystretch}{2em}
\multlinegap0pt
\usepackage{amssymb}
\usepackage{xcolor}
\usepackage{cancel}
\usepackage[normalem]{ulem}
\usepackage[shortlabels]{enumitem}
\usepackage{tikz}
\newtheorem{lemma}{Lemma}
\newcommand{\R}{\mathbb{R}}
\newcommand{\md}{\mathrm{d}}
\title{Hyperuniformity and optimal transport of point processes}
\author{Rapha\"el Lachi\`eze-Rey, D. Yogeshwaran}
\date{Source: arXiv:2402.13705v4 (CC BY 4.0)}
\begin{document}
\maketitle

\noindent\emph{Source and licence.} Hyperuniformity and optimal transport of point processes. Version arXiv:2402.13705v4 (CC BY 4.0), distributed by arXiv under the Creative Commons Attribution 4.0 International licence, http://creativecommons.org/licenses/by/4.0/, as recorded in the arXiv licence metadata for this exact version and shown on the abstract page https://arxiv.org/abs/2402.13705v4. Full licence notice: \texttt{licenses/arxiv-2402.13705-CC-BY-4.0.txt}.\par\smallskip

\noindent\textit{Editorial scope.} Counted unit: the article's lemma lm:gamma\_n (source0.tex lines 1256--1268). The article's complete original proof of it is delivered with it (source0.tex lines 1271--1299, one \texttt{proof} environment of 29 lines). No mathematics is written, repaired or re-derived: the statement and the proof are the article's own lines, in the article's own notation, and no retained line is substituted or re-flowed.  No further source-local context is needed: every cross-reference of every retained line resolves inside the two delivered spans.  Nothing was omitted from any retained span, so the package carries no omission record.  No retained line contains a citation, so the wrapper carries no bibliography and no bibliography entry.  No retained line contains a figure environment, an image-inclusion command or drawing code, so no image is delivered and the package declares no asset dependency.  Editorial formatting only, in addition: an AMS \texttt{amsart} preamble that restates the article's own declarations and macros used by the retained lines, namely the environments \texttt{lemma} on separate counters, together with the standard packages those lines need.  The retained environments print in this document as Lemma 1 in order of appearance, whereas the article prints the same items with the numbering of its own full document.  The complete native source of the article as distributed in the arXiv e-print for this exact version is bundled unchanged as \texttt{sources/154910/source0.tex}.
\begin{lemma}
\label{lm:gamma_n}
We have for all $x \in \R^d, n \in \mathbb{N}$,
\begin{align}
\label{eq:bd-gamma_n}
0<c\leqslant \frac{ \gamma _{n}(x) }{n\min(1,\|x\|n^{-1/d})}\leqslant C<\infty
\end{align}
and hence
\begin{align*}
 | \int_{}\gamma _{n}(x)\beta (\md x) | \leqslant   c \, n \, \varepsilon (n^{-1/d}).
\end{align*}
The reverse inequality holds (up to a constant) if $ \beta $ has a constant sign.
\end{lemma}

\begin{proof}  {We shall prove \eqref{eq:bd-gamma_n} with $\ell_{\infty}$ norm and by equivalence of norms in $\R^d$, the statement holds also for Euclidean norm. Set $ \|x\| = \max_{i} | x_{i} |$ and $x_+  := \max \{x,0\}$ for $x \in \R$.}
{We have the explicit expressions
%
\begin{align*}
\delta _{n}(x)= \prod_{i= 1}^{d}(2\pi n^{1/d}- | x_{i} | )_{+}, \, \, \text{and} \, \, \gamma _{n}(x)&
 = \sum_{j =  1}^{d}(2\pi n^{ 1/d})^{j- 1} | x_{j} |  \prod_{k = j +  1}^{d}(2\pi n^{ 1/d}- | x_{k} | )_+,
\end{align*}
%
where the latter uses the following telescopic formula: for real numbers $ a_{i},b_{i}$, with the convention $ b_{d +  1} = a_{0} =  1, $ it holds that
%
\begin{align*}  \prod_{i}a_{i}-  \prod_{i}b_{i} = &\sum_{j =   1}^{d}(a_{ 1}\dots a_{j}b_{j +  1 }\dots b_{d}-a_{ 1}\dots a_{j- 1 }b_{j}\dots b_{d}),
\end{align*}
%
}If $ \|x\|\geqslant 2\pi n^{ 1/d}$, $ \gamma  _{n}(x) = (2\pi )^{d}n$, and the result holds. We have the general upper bound if $ \|x\|<2\pi n^{ 1/d}$ as then $ \gamma _{n}(x)\leqslant {d}\|x\|(2\pi n)^{\frac{ d- 1}{d}}\leqslant n$. We need to prove a lower bound in this case.

{Since $ \|x\| < 2\pi n^{ 1/d}$, there is} $1 \leqslant i_{0} \leqslant d$ such that
%
\begin{align*}
|x_{i_{0}}|  = \|x\| <2\pi n^{ 1/d}.
\end{align*}
%
Now if $ \|x\|>\pi n^{ 1/d}$ then $| \delta _{n}(x)| < | \Lambda _{n} | /2$, and $ \gamma _{n}(x)\geqslant cn,$ hence the lower bound follows in this case.
Else $ \|x\|< \pi n^{ 1/d}$ and then $ 2\pi n^{ 1/d}- | x_{k} | \geqslant  \pi n^{ 1/d}$ for all $k$, and so one can lower bound by the single term $ j = i_{0}:$
%
\begin{align*}\gamma _{n}(x)\geqslant cn^{\frac{ d- 1}{d}}|x_{i_{0}}| = c\|x\|n^{ 1- 1/d}.
\end{align*}
%
{Using the upper bound for $\gamma_n$ and the definition of $\varepsilon$, one can immediately bound $\int_{}\gamma _{n}(x)|\beta|(\md x)$ and in case $\beta$ has constant sign, we derive the lower bound using the lower bound for $\gamma_n$.}
\end{proof}

\end{document}
